% Integration fragment; theorem/proof environments must be defined by the host.
% All variables in the learning and safety subsections are dimensionless.

\subsection{Rod mechanics and a variationally consistent reduction}
Following the spatial rod formulation \citep{simo,antman}, let $s\in[0,L]$ denote reference arclength, $r(s,t)\in\mathbb R^3$
the centreline and $D(s,t)\in SO(3)$ its director frame. Define
\begin{equation}
 \Gamma=D^{\mathsf T}r_s,\quad
 D_s=D\widehat\kappa,\quad D_t=D\widehat\omega_b.
\end{equation}
The hat map is defined by $\widehat a b=a\times b$.
The extensible, unshearable Kirchhoff specialization has
$\Gamma=(0,0,1+e)^{\mathsf T}$; inextensibility additionally sets $e=0$.
Transverse internal forces are then constraint reactions. They must not
be set to zero by substituting zero shear into a shearable constitutive law.
In spatial components the balances are
\begin{align}
 \rho A r_{tt}&=n_s+f_{\rm mag}+f_{\rm h}+f_{\rm c}+\rho A g,\\
 \partial_t(D\mathcal J_b\omega_b)&=m_s+r_s\times n+\tau_{\rm ext}.
\end{align}
For constant body inertia, the left side of the second equation is
$D(\mathcal J_b\dot\omega_b+\omega_b\times\mathcal J_b\omega_b)$.
With stored energy density $W(\Gamma,\kappa)$, unconstrained components
of $D^{\mathsf T}n$ and $D^{\mathsf T}m$ are its strain derivatives.
A fixed base specifies position and orientation; a free tip specifies
the appropriate force and moment resultants, including any concentrated
magnetic load.

For a body-fixed magnetic moment density $\mu_b(s)$ and a static applied
field $B_0(r)$, neglecting magnetic interactions between distinct
material elements gives the applied-field potential
\begin{equation}
 V_{\rm mag}=-\int_0^L(D\mu_b)\cdot B_0(r)\,ds.
\end{equation}
Its translational and rotational variations give
$f_{\rm mag}=\nabla_r(\mu\cdot B_0)$ at fixed orientation and
$\tau_{\rm mag}=\mu\times B_0$, where $\mu=D\mu_b$.
Thus static magnetic loading belongs to the conservative energy.
Field modulation supplies explicit time-dependent work or an input port.
The permanent-dipole load is the leading term of the Maxwell-stress
resultant derived below; it is counted once in the force balance.

For a complete localized magnetic element, the magnetostatic force
can be evaluated on a surrounding vacuum control surface $\Sigma$:
\begin{equation}
 F_{\rm mag}=\int_\Sigma T_Mn_\Sigma\dd A,\qquad
 T_M=\mu_0^{-1}(BB^{\mathsf T}-\tfrac12|B|^2I).
\end{equation}
Here $B$ is the total field, including the element's self-field.
The associated moment integrates the moment arm crossed with
$T_Mn_\Sigma$. To recover the permanent-dipole approximation, represent
the element by its complete localized steady current $j(\xi)$,
including bound currents, with $\xi=x-r$. The prescribed magnetization
is unaffected by the applied field $B_{\rm ext}$, whose sources lie
outside $\Sigma$. Current closure gives $\nabla\cdot j=0$ and no flux
through the boundary. The identity $\nabla\cdot T_M=j\times B$ and
symmetry of $T_M$ convert the stress resultants into Lorentz force and
moment. The self-field has zero net force and moment, as seen by taking
its stress surface to infinity; the applied-field-only stress also has
zero resultant in this source-free volume. The remaining interaction is
\begin{equation}
\begin{aligned}
 m_{\rm d}&=\tfrac12\int\xi\times j(\xi)\,\dd V,\\
 F_{\rm mag}&=\int j(\xi)\times B_{\rm ext}(r+\xi)\,\dd V\\
 &=\nabla_r\!\left(m_{\rm d}\cdot B_{\rm ext}(r)\right)+R_F,\\
 \tau_{\rm mag}&=\int\xi\times[j(\xi)\times B_{\rm ext}(r+\xi)]\,\dd V\\
 &=m_{\rm d}\times B_{\rm ext}(r)+R_\tau .
\end{aligned}
\label{eq:stress-dipole-reduction}
\end{equation}
The gradient holds the moment orientation fixed.
Integration by parts gives $\int j\,\dd V=0$ and
$\int\xi_i j_j\,\dd V=\epsilon_{ijk}(m_{\rm d})_k$.
Taylor expansion of $B_{\rm ext}$ and $\nabla\cdot B_{\rm ext}=0$
then give \eqref{eq:stress-dipole-reduction}.
For $J_2=\int|\xi|^2|j(\xi)|\,\dd V$, the remainders obey
$|R_F|\le\tfrac12J_2\sup\|D^2B_{\rm ext}\|$ and
$|R_\tau|\le J_2\sup\|DB_{\rm ext}\|$, with operator norms taken over
the element and its Taylor segments. Surface-current contributions are
included in these integrals when present. Neglecting the higher
multipoles recovers the permanent-dipole force and torque used in
magnetic robotics \citep[Section~4.1, Eqs.~(12) and~(16)]{abbott}.

A modified Morison-type local closure, stated per unit reference length,
is
\begin{align}
 f_{\rm h}={}&m_f a_f+m_a(a_f-r_{tt})
             -c_1v_r-c_2|v_r|v_r,\\
 v_r={}&r_t-v_f,\qquad c_1,c_2\ge0.
\end{align}
Here $m_f=\rho_f A_f$ and $m_a$ are displaced-fluid and added mass
per unit length; $a_f$ is the prescribed fluid acceleration.
This is a reduced closure requiring calibration, not a solved
Navier--Stokes interaction. In still fluid its drag power is
$-c_1|r_t|^2-c_2|r_t|^3$. In moving fluid it is dissipative relative to
the fluid velocity and can do positive work on the rod. A constant
added mass may be moved to the inertial side exactly once.

For a stationary wall, define penetration $\delta(r)$ positive inside
the wall and $n_c=\nabla\delta$. One nonadhesive penalty model is
\begin{align}
 V_c&=\frac{k_c}{m+1}[\delta]_+^{m+1},\\
 f_c&=-k_c[\delta]_+^m n_c
       -c_c\mathbf1_{\{\delta>0\}}[\dot\delta]_+n_c.
\end{align}
The elastic term is included through $V_c$ and the second term dissipates
power during compression. Friction can be added as a dissipative port.
This switching law is used for force accounting outside the contact-free
KAM domain, not as an analytic perturbation through an impact.

Choose a configuration map $(r(s,q),D(s,q))$ that satisfies the essential
boundary conditions and the rod constraints. A linear centreline
ansatz alone is insufficient to specify rotational inertia. Define
\begin{equation}
 r_i=\partial_{q_i}r,\qquad
 a_i=\operatorname{vee}(D^{\mathsf T}\partial_{q_i}D).
\end{equation}
Then $\omega_b=\sum_i a_i\dot q_i$ and the reduced kinetic energy is
\begin{align}
 T_N&=\tfrac12\dot q^{\mathsf T}M(q)\dot q,\\
 M_{ij}(q)&=\int_0^L\left(\rho A r_i\cdot r_j+
                 a_i^{\mathsf T}\mathcal J_ba_j\right)ds.
\end{align}
An added-fluid-mass contribution may be included when its consistent
kinetic metric and associated inertial terms are specified. Otherwise
hydrodynamic inertia remains in the generalized force balance and is
not also included in $M$. Galerkin projection of the virtual-work
identity gives
\begin{equation}
 M(q)\ddot q+C(q,\dot q)\dot q+\nabla V_N(q)=Q_{\rm nc},
\end{equation}
\begin{equation}
 Q_i=\int_0^L(f\cdot r_i+\tau\cdot Da_i)ds+Q_i^\partial.
\end{equation}
Here $C\dot q$ denotes the Christoffel contribution of $M$, and $V_N$
contains elastic, static magnetic and retained gravitational potentials.
Microgravity reduces the last contribution; it does not eliminate
fluid drag or establish a particular blood-flow profile.
Relative-flow drag, moving-base work and contact losses are retained in
$Q_{\rm nc}$. A stationary-fluid quadratic drag, for example, has
nonpositive power; a moving fluid may inject energy. Contact switching
is excluded from the analytic neighbourhood used for KAM analysis.

For $M(q)\succ0$, the Legendre transform yields
\begin{equation}
 p=M(q)\dot q,\qquad
 H_N(q,p)=\tfrac12p^{\mathsf T}M(q)^{-1}p+V_N(q).
\end{equation}
In particular,
$\dot p_i=-\tfrac12p^{\mathsf T}\partial_{q_i}(M^{-1})p
-\partial_{q_i}V_N+(Q_{\rm nc})_i$.
The configuration dependence of inertia is retained. Potential forces
appear either in $V_N$ or in $Q_{\rm nc}$, never both.

\subsection{Complex energy--dissipation representation with an explicit decoder}
Energy and Rayleigh dissipation have different physical units. Let $E_0$
and $t_0$ be reference energy and time, and choose diagonal position and
momentum scales $Q_0,P_0$ with $Q_0^{\mathsf T}P_0=E_0t_0I$.
Set $q=Q_0x$, $p=P_0y$, $\tau=t/t_0$ and
\begin{align}
 h(x,y)&=H_N(Q_0x,P_0y)/E_0,\\
 d(x,v)&=\frac{t_0}{E_0}\mathcal D_N(Q_0x,Q_0v/t_0).
\end{align}
Here $v=dx/d\tau$ and $\mathcal D_N$ is a Rayleigh dissipation potential.
The two entries of $h+id$ are now dimensionless. A complex network
provides coupled features and two physically decoded real readouts,
\begin{align}
 \mathcal H_\theta(x,y,v)&=h_\theta(x,y)+i d_\theta(x,v),\\
 h_\theta&=h_{\rm prior}+h_{{\rm res},\theta}.
\end{align}
One admissible dissipative head is
\begin{equation}
 d_\theta(x,v)=\tfrac12v^{\mathsf T}D_\theta(x)v,
 \qquad D_\theta(x)=L_\theta(x)L_\theta(x)^{\mathsf T}\succeq0.
\end{equation}
More general heads may be used if $d_\theta(x,\cdot)$ is differentiable,
convex and minimized at zero. The decoder is
\begin{align}
 v&=\partial_y h_\theta,\qquad x'=v,\\
 y'&=-\partial_x h_\theta-\partial_vd_\theta(x,v)+B(x)u.
 \label{eq:fragment-decoder}
\end{align}
The derivative of $d_\theta$ is taken with respect to its independent
velocity argument before substituting $v=\partial_yh_\theta$.
It is not a total derivative through this substitution.

\begin{proposition}[Power-consistent complex representation]
For a differentiable $h_\theta$, the decoded system satisfies
\begin{equation}
 \frac{d h_\theta}{d\tau}
 =-v^{\mathsf T}\partial_vd_\theta(x,v)
   +u^{\mathsf T}B(x)^{\mathsf T}v+\partial_\tau h_\theta.
 \label{eq:fragment-power}
\end{equation}
For the quadratic head, the first term equals
$-v^{\mathsf T}D_\theta(x)v\le0$. For the convex head it is also
nonpositive.
\end{proposition}
\begin{proof}
Apply the chain rule to $h_\theta$ and substitute
\eqref{eq:fragment-decoder}. The two canonical terms cancel.
Convexity and the minimum at zero imply
$v^{\mathsf T}\partial_vd_\theta(x,v)
\ge d_\theta(x,v)-d_\theta(x,0)\ge0$.
\end{proof}

With $z=(x+iy)/\sqrt2$, the canonical convention is
$\{z_j,\bar z_k\}=-i\delta_{jk}$ and
\begin{equation}
 z'=-i\partial_{\bar z}h_\theta
       -\frac{i}{\sqrt2}\partial_vd_\theta
       +\frac{i}{\sqrt2}B(x)u.
\end{equation}
This is a unified complex representation of conservative and
nonconservative dynamics with a specified physical decoder. The rule
$z'=-i\partial_{\bar z}(h_\theta+i d_\theta)$ would generally produce
a different vector field and has no automatic dissipativity guarantee.
Real analyticity in $(x,y)$, and a controlled analytic extension to a
complex neighbourhood, are the relevant requirements for analytic KAM.
Holomorphicity in $z$ alone is not required of a nonconstant real energy
$h_\theta(z,\bar z)$. In particular, ReLU or modReLU features require a
separate smooth approximation or a theorem of appropriate regularity;
their ordinary piecewise-smooth implementation is not analytic.


The dimensional specialization below uses explicit SI units. The learning
and KAM statements use dimensionless canonical variables, with their
energy and time scales declared before fitting.
